\documentclass[12pt]{article}
\usepackage{amsmath}
\usepackage{amsfonts}
\usepackage{float}
\usepackage{fancyhdr}
\usepackage{graphicx}
\usepackage{setspace}
\usepackage[colorlinks=true,linkcolor=blue, citecolor=red]{hyperref}
\usepackage{url}
\usepackage[top=.75in, left=.75in, right=.75in, bottom=1in]{geometry}

% For algorithm
\usepackage{algorithm}
\usepackage{algpseudocode}

% ============ CODE ============
\usepackage{listings}
\usepackage{xcolor}
\definecolor{codegreen}{rgb}{0,0.6,0}
\definecolor{codegray}{rgb}{0.5,0.5,0.5}
\definecolor{codepurple}{rgb}{0.58,0,0.82}
\definecolor{backcolour}{rgb}{0.95,0.95,0.92}

% Styling for the code.
\lstdefinestyle{mystyle}{
    backgroundcolor=\color{backcolour},   
    commentstyle=\color{codegreen},
    keywordstyle=\color{magenta},
    numberstyle=\tiny\color{codegray},
    stringstyle=\color{codepurple},
    basicstyle=\ttfamily\footnotesize,
    breakatwhitespace=false,         
    breaklines=true,                 
    captionpos=b,                    
    keepspaces=true,                 
    numbers=left,                    
    numbersep=5pt,                  
    showspaces=false,                
    showstringspaces=false,
    showtabs=false,                  
    tabsize=2
}
\lstset{style=mystyle}

% Table packages and Unicode / Font support
\usepackage{array}
\usepackage{longtable}
\usepackage{enumitem}
\usepackage{booktabs}
\usepackage[utf8]{inputenc}
\DeclareUnicodeCharacter{2264}{\ensuremath{\le}}
\DeclareUnicodeCharacter{2194}{\ensuremath{\leftrightarrow}}
\usepackage[T1]{fontenc}
\usepackage{lmodern}
\graphicspath{{../../../../templates/}{../../prototype/}{figures/}{./}}
\setlength{\emergencystretch}{3em}
\setlist[enumerate]{label=\arabic*.,leftmargin=*}
\setlist[itemize]{leftmargin=*}

% ============ HEADER AND FOOTER ============
\setlength{\headheight}{29.43912pt}

\pagestyle{fancy}
\fancyfoot{}
\fancyfoot[R]{Page \thepage}

\setlength{\parindent}{0pt}
\renewcommand{\baselinestretch}{1.5}

\newcommand{\coursename}{Introduction of Software Engineering}
\newcommand{\reportname}{Studitopia}
\newcommand{\reporttitle}{\href{https://github.com/RoundFloor/Studitopia.git}{Studitopia}}

\newcommand{\studentname}{%
Nguyen Thanh An (24125024)\\
Nguyen Chi Bao (24125025)\\
Tran Nguyen Khanh Binh (24125026)\\
Tran Canh Anh Tuan (24125048)\\
Nguyen Le Quoc Huy (24125100)
}
\newcommand{\teachername}{Ho Tuan Thanh, Nguyen Thi Minh Tuyen}

% Header
\lhead{\reporttitle}
\rhead{
University of Science - VNUHCM\\
\coursename
}

% Footer
\newcommand{\leftfooter}{}
\lfoot{\leftfooter}

% ============ DOCUMENT ============
\begin{document}

\pagenumbering{roman}
\begin{titlepage}
\newcommand{\HRule}{\rule{\linewidth}{0.5mm}}
\centering

\textsc{\LARGE VNU-HCM}\\[0.5cm]
\textsc{\Large University of Science}\\[0.5cm]
\textsc{\large Faculty of Information Technology}\\[0.5cm]
\textsc{Computer Science}\\[0.5cm]

\HRule \\[0.4cm]
{ 
\huge{\bfseries{\reporttitle}}\\[0.5cm]
\large{\bfseries{UC-02 --- Log In}}
}\\[0.4cm]
\HRule \\[0.5cm]

\textbf{\large Subject: \coursename}\\[0.5cm]

\begin{minipage}[t]{0.55\textwidth}
\begin{flushleft} \large
\emph{Students:}\\
\studentname
\end{flushleft}
\end{minipage}
~
\begin{minipage}[t]{0.4\textwidth}
\begin{flushright} \large
\emph{Supervising Teachers:} \\
\teachername
\end{flushright}
\end{minipage}\\[1cm]

{\large \today}\\[1cm]

\includegraphics[scale=.20]{img/hcmus-logo.png}\\[1cm] 

\vfill
\end{titlepage}	

\tableofcontents
\pagebreak

\pagenumbering{arabic}
\setcounter{page}{1}

\hypertarget{sec:uc-02-log-in}{}
\addcontentsline{toc}{section}{UC-02 --- Log In}
\section*{UC-02 --- Log In}
\par\medskip

\textbf{Version:} \texttt{2.0}
\par\medskip

\textbf{Performed by:} [Nguyen Le Quoc Huy] | \textbf{Edited by:} [Nguyen Thanh An]
\par\medskip

\begingroup
\small
\renewcommand{\baselinestretch}{1.0}\selectfont
\renewcommand{\arraystretch}{1.15}
\setlength{\tabcolsep}{6pt}
\begin{longtable}{|>{\raggedright\arraybackslash}p{0.09\textwidth}|>{\raggedright\arraybackslash}p{0.84\textwidth}|}
\hline
\textbf{Field} & \textbf{Value} \\ \hline
\endhead
Use Case ID & UC-02 \\ \hline
Name & Log In \\ \hline
Primary Actor & Registered University Student \\ \hline
Supporting Actors & Google OAuth 2.0 \\ \hline
Feature Group & Authentication \\ \hline
Priority & High \\ \hline
Related User Stories & US-02 \\ \hline
Related Requirements & US-02, US-04; FR-02.1 to FR-02.5, FR-04.1 to FR-04.4; NFR-SEC-1, NFR-SEC-3, NFR-SEC-6, NFR-SEC-7, NFR-SEC-8 \\ \hline
Source Documents & \texttt{use\_cases\_model/UC-02-log-in.md}; \texttt{use\_cases\_description.md} (UC-02); \texttt{use\_cases\_model.md}; \texttt{use\_case\_diagrams.md}; \texttt{PA3-Group05/use\_case.md} (§2 UC-02); \texttt{overview.md}; \texttt{functional\_requirement.md}; \texttt{non\_functional\_requirement.md}; \texttt{userstory.md}; \texttt{feature\_flow.md}; \texttt{screen\_flow.md}; \texttt{API.md}; \texttt{response\_code.md}; \texttt{settings.md} \\ \hline
\end{longtable}
\endgroup
\par\medskip

\hypertarget{sec:part-a-use-case-model}{}
\addcontentsline{toc}{subsection}{Part A --- Use Case Model}
\subsection*{Part A --- Use Case Model}
\par\medskip

\hypertarget{sec:a1-focused-use-case-diagram}{}
\addcontentsline{toc}{subsubsection}{A1. Focused Use Case Diagram}
\subsubsection*{A1. Focused Use Case Diagram}
\par\medskip

\begin{center}
\includegraphics[width=0.85\textwidth,height=0.75\textheight,keepaspectratio]{"figures/UC-02-log-in_1.png"}
\end{center}
\par\medskip

\textbf{Diagram notation.} Person shape = human actor; cloud = supporting system; blue = selected base use case; purple = included; green = extending; solid line = association; dashed arrow = include or extend; labelled solid arrow = generalization.
\par\medskip

\hypertarget{sec:a2-actor-register}{}
\addcontentsline{toc}{subsubsection}{A2. Actor Register}
\subsubsection*{A2. Actor Register}
\par\medskip

\begingroup
\small
\renewcommand{\baselinestretch}{1.0}\selectfont
\renewcommand{\arraystretch}{1.15}
\setlength{\tabcolsep}{6pt}
\begin{longtable}{|>{\raggedright\arraybackslash}p{0.16\textwidth}|>{\raggedright\arraybackslash}p{0.10\textwidth}|>{\raggedright\arraybackslash}p{0.66\textwidth}|}
\hline
\textbf{Actor} & \textbf{Type} & \textbf{Description} \\ \hline
\endhead
Registered University Student & Primary & A student with an existing account who logs in to access protected features. \\ \hline
University Student & General human actor & Parent of Registered University Student; inherits general student associations to UC-03, UC-04, UC-05, UC-06, and UC-07. \\ \hline
Google OAuth 2.0 & Supporting System & Provides an optional external authentication and registration path; supports E-02 Log In with Google. \\ \hline
\end{longtable}
\endgroup
\par\medskip

\hypertarget{sec:a3-relationships}{}
\addcontentsline{toc}{subsubsection}{A3. Relationships}
\subsubsection*{A3. Relationships}
\par\medskip

\hypertarget{sec:a3-1-associations}{}
\paragraph*{A3.1 Associations}
\par\medskip

\begingroup
\small
\renewcommand{\baselinestretch}{1.0}\selectfont
\renewcommand{\arraystretch}{1.15}
\setlength{\tabcolsep}{6pt}
\begin{longtable}{|>{\raggedright\arraybackslash}p{0.21\textwidth}|>{\raggedright\arraybackslash}p{0.17\textwidth}|>{\raggedright\arraybackslash}p{0.10\textwidth}|>{\raggedright\arraybackslash}p{0.44\textwidth}|}
\hline
\textbf{Actor} & \textbf{Use Case} & \textbf{Type} & \textbf{Description} \\ \hline
\endhead
Registered University Student & UC-02 Log In & Primary & Authenticates with credentials and reaches protected features. \\ \hline
Google OAuth 2.0 & E-02 Log In with Google & Supporting & Authenticates the optional Google login path. \\ \hline
\end{longtable}
\endgroup
\par\medskip

\hypertarget{sec:a3-2-generalization-and-specialization}{}
\paragraph*{A3.2 Generalization and Specialization}
\par\medskip

\begingroup
\small
\renewcommand{\baselinestretch}{1.0}\selectfont
\renewcommand{\arraystretch}{1.15}
\setlength{\tabcolsep}{6pt}
\begin{longtable}{|>{\raggedright\arraybackslash}p{0.21\textwidth}|>{\raggedright\arraybackslash}p{0.13\textwidth}|>{\raggedright\arraybackslash}p{0.15\textwidth}|>{\raggedright\arraybackslash}p{0.42\textwidth}|}
\hline
\textbf{Child} & \textbf{Parent} & \textbf{Type} & \textbf{Inherited Behavior or Association} \\ \hline
\endhead
Registered University Student & University Student & Actor generalization & Inherits the general student associations and adds UC-02. \\ \hline
\end{longtable}
\endgroup
\par\medskip

\hypertarget{sec:a3-3-include}{}
\paragraph*{A3.3 Include}
\par\medskip

\begingroup
\small
\renewcommand{\baselinestretch}{1.0}\selectfont
\renewcommand{\arraystretch}{1.15}
\setlength{\tabcolsep}{6pt}
\begin{longtable}{|>{\raggedright\arraybackslash}p{0.10\textwidth}|>{\raggedright\arraybackslash}p{0.26\textwidth}|>{\raggedright\arraybackslash}p{0.56\textwidth}|}
\hline
\textbf{Base Use Case} & \textbf{Included Use Case} & \textbf{Meaning} \\ \hline
\endhead
UC-02 Log In & I-03 Create Authenticated Session & Create the authenticated session and return the student's access context. \\ \hline
\end{longtable}
\endgroup
\par\medskip

\hypertarget{sec:a3-4-extend}{}
\paragraph*{A3.4 Extend}
\par\medskip

\begingroup
\small
\renewcommand{\baselinestretch}{1.0}\selectfont
\renewcommand{\arraystretch}{1.15}
\setlength{\tabcolsep}{6pt}
\begin{longtable}{|>{\raggedright\arraybackslash}p{0.16\textwidth}|>{\raggedright\arraybackslash}p{0.10\textwidth}|>{\raggedright\arraybackslash}p{0.38\textwidth}|>{\raggedright\arraybackslash}p{0.28\textwidth}|}
\hline
\textbf{Extending Use Case} & \textbf{Base Use Case} & \textbf{Condition} & \textbf{Added Behavior} \\ \hline
\endhead
E-02 Log In with Google & UC-02 Log In & The student selects Google login instead of credentials. & Add the Google OAuth authentication path. \\ \hline
\end{longtable}
\endgroup
\par\medskip

\hypertarget{sec:a4-model-notes}{}
\addcontentsline{toc}{subsubsection}{A4. Model Notes}
\subsubsection*{A4. Model Notes}
\par\medskip

\begin{enumerate}
\item Successful login always includes I-03 because a valid authenticated session is required before protected features are available.
\item Google login is an optional extension of UC-02. The standard username-or-email and password path remains the base behavior.
\item Registered University Student inherits all general University Student associations in addition to the direct UC-02 association.
\end{enumerate}
\par\medskip

\hypertarget{sec:part-b-use-case-specification}{}
\addcontentsline{toc}{subsection}{Part B --- Use Case Specification}
\subsection*{Part B --- Use Case Specification}
\par\medskip

\hypertarget{sec:b1-brief-description}{}
\addcontentsline{toc}{subsubsection}{B1. Brief Description}
\subsubsection*{B1. Brief Description}
\par\medskip

A registered university student authenticates and regains access to their profile, university forum, posts, and learning activities. Login supports a registered username or university email with a password, and an optional Google OAuth path.
\par\medskip

\hypertarget{sec:b2-preconditions}{}
\addcontentsline{toc}{subsubsection}{B2. Preconditions}
\subsubsection*{B2. Preconditions}
\par\medskip

\begin{enumerate}
\item The student has a registered and activated account.
\item The student is not already authenticated, or chooses to start a new authenticated session.
\item The authentication service and, when used, Google OAuth 2.0 are available.
\end{enumerate}
\par\medskip

\hypertarget{sec:b3-postconditions}{}
\addcontentsline{toc}{subsubsection}{B3. Postconditions}
\subsubsection*{B3. Postconditions}
\par\medskip

\begin{enumerate}
\item \textbf{Success:} A valid authenticated session is created and the student reaches the Home page or the originally requested protected page.
\item \textbf{Success:} The student's default university forum context is available after login.
\item \textbf{Failure:} No authenticated session is created, and the student remains on the login page with a non-disclosing error message.
\end{enumerate}
\par\medskip

\hypertarget{sec:b4-main-success-scenario-basic-flow}{}
\addcontentsline{toc}{subsubsection}{B4. Main Success Scenario (Basic Flow)}
\subsubsection*{B4. Main Success Scenario (Basic Flow)}
\par\medskip

\begin{enumerate}
\item The student opens the login page.
\item The system displays the username-or-email field, password field, CAPTCHA, \textbf{Remember Me} option where supported, \textbf{Forget Password}, \textbf{Register}, and \textbf{Continue with Google} options.
\item The student enters a registered username or university email address and password, then completes the CAPTCHA.
\item The frontend performs basic input validation and submits the login request.
\item The backend verifies the CAPTCHA, credentials, account activation status, and any account restrictions.
\item The backend creates an authenticated session and returns the access result.
\item The frontend stores the session according to the \textbf{Remember Me} choice, loads the student's profile and forum context, and redirects to Home or the intended protected page.
\item The system displays the Home page with the Forum entry and the student's university forum as the default context.
\end{enumerate}
\par\medskip

\hypertarget{sec:b5-alternative-flows}{}
\addcontentsline{toc}{subsubsection}{B5. Alternative Flows}
\subsubsection*{B5. Alternative Flows}
\par\medskip

Unless otherwise stated, an alternative flow resumes at the next step of the basic flow.
\par\medskip

\hypertarget{sec:a1-google-login}{}
\paragraph*{A1. Google login}
\par\medskip

At basic flow step 2, the student selects \textbf{Continue with Google}. The system completes Google OAuth, validates the response, and either creates the authenticated session or displays an appropriate error.
\par\medskip

\hypertarget{sec:a2-remember-me-selected}{}
\paragraph*{A2. Remember Me selected}
\par\medskip

At basic flow step 7, the system maintains the session across browser restarts according to the selected option and the platform session policy.
\par\medskip

\hypertarget{sec:a3-intended-page-available}{}
\paragraph*{A3. Intended page available}
\par\medskip

If the student was redirected from a protected page, the system returns the student to that page after successful authentication instead of Home.
\par\medskip

\hypertarget{sec:a4-unverified-account}{}
\paragraph*{A4. Unverified account}
\par\medskip

If the backend finds that the account is not verified, login is denied and the student is directed to complete OTP verification.
\par\medskip

\hypertarget{sec:b6-exception-flows}{}
\addcontentsline{toc}{subsubsection}{B6. Exception Flows}
\subsubsection*{B6. Exception Flows}
\par\medskip

\hypertarget{sec:e1-invalid-credentials}{}
\paragraph*{E1. Invalid credentials}
\par\medskip

The system displays a clear error without revealing whether the username, email, or password was incorrect.
\par\medskip

\hypertarget{sec:e2-captcha-failure}{}
\paragraph*{E2. CAPTCHA failure}
\par\medskip

The system rejects the request, displays an error, and allows the student to retry.
\par\medskip

\hypertarget{sec:e3-unavailable-or-restricted-account}{}
\paragraph*{E3. Unavailable or restricted account}
\par\medskip

The system denies authentication and displays an appropriate message without exposing sensitive account details.
\par\medskip

\hypertarget{sec:e4-authentication-service-failure}{}
\paragraph*{E4. Authentication service failure}
\par\medskip

No session is created. The system displays a temporary error and allows the student to retry.
\par\medskip

\hypertarget{sec:e5-session-expires-while-accessing-a-protected-page}{}
\paragraph*{E5. Session expires while accessing a protected page}
\par\medskip

The system denies access, informs the student that the session expired, and requires login again.
\par\medskip

\hypertarget{sec:b7-business-rules}{}
\addcontentsline{toc}{subsubsection}{B7. Business Rules}
\subsubsection*{B7. Business Rules}
\par\medskip

\begin{enumerate}
\item Login is accepted with either a registered username or a registered university email address.
\item Invalid-credential messages must not reveal which credential field was incorrect.
\item CAPTCHA must be successfully verified before login processing.
\item Google login is an optional authentication path for registered users.
\item Protected pages and APIs require a valid authenticated session.
\item After successful login, the student's profile, forum, posts, and learning activities are accessible according to their permissions.
\item The session is subject to the platform's inactivity-expiry policy.
\end{enumerate}
\par\medskip

\hypertarget{sec:b8-special-requirements}{}
\addcontentsline{toc}{subsubsection}{B8. Special Requirements}
\subsubsection*{B8. Special Requirements}
\par\medskip

\begingroup
\small
\renewcommand{\baselinestretch}{1.0}\selectfont
\renewcommand{\arraystretch}{1.15}
\setlength{\tabcolsep}{6pt}
\begin{longtable}{|>{\raggedright\arraybackslash}p{0.10\textwidth}|>{\raggedright\arraybackslash}p{0.83\textwidth}|}
\hline
\textbf{ID} & \textbf{Requirement} \\ \hline
\endhead
NFR-SEC-1 & Authentication is required before accessing protected user resources. \\ \hline
NFR-SEC-3 & Sensitive communication must be encrypted using TLS 1.2 or higher. \\ \hline
NFR-SEC-6 & Server-side authorization checks must prevent unauthorized access. \\ \hline
NFR-SEC-7 & Inactive authentication sessions must expire automatically after 30 minutes. \\ \hline
NFR-SEC-8 & The system must protect against SQL injection, XSS, CSRF, and other common web vulnerabilities. \\ \hline
\end{longtable}
\endgroup
\par\medskip

\hypertarget{sec:b9-extension-points}{}
\addcontentsline{toc}{subsubsection}{B9. Extension Points}
\subsubsection*{B9. Extension Points}
\par\medskip

\begingroup
\small
\renewcommand{\baselinestretch}{1.0}\selectfont
\renewcommand{\arraystretch}{1.15}
\setlength{\tabcolsep}{6pt}
\begin{longtable}{|>{\raggedright\arraybackslash}p{0.39\textwidth}|>{\raggedright\arraybackslash}p{0.54\textwidth}|}
\hline
\textbf{Extension Point} & \textbf{Location in Flow of Events} \\ \hline
\endhead
Google login & Basic flow step 2; alternative flow A1 \\ \hline
Remember Me & Basic flow step 7; alternative flow A2 \\ \hline
Return to intended page & Basic flow step 8; alternative flow A3 \\ \hline
Unverified-account recovery & Basic flow step 5; alternative flow A4 \\ \hline
Session expiry & Protected-page access; exception E5 \\ \hline
\end{longtable}
\endgroup
\par\medskip

\hypertarget{sec:b10-interface-and-data-notes}{}
\addcontentsline{toc}{subsubsection}{B10. Interface and Data Notes}
\subsubsection*{B10. Interface and Data Notes}
\par\medskip

\begin{itemize}
\item Screen/flow path: \texttt{Start → IsLogin (False) → Authentication → Login}; successful \texttt{Login → Home}; \texttt{Login → Register}, \texttt{Login → Recover}, \texttt{Login → Home}; \texttt{Home → Settings → Logout → Login}; feature flow \texttt{Authentication → Login → Home}.
\item Login UI fields and actions: Username / Email, Password, CAPTCHA, \textbf{Remember Me} where supported, \textbf{Forget Password}, \textbf{Register}, and \textbf{Continue with Google}.
\item \texttt{POST /auth/login} query \texttt{type: 'token' | 'username' | 'email'}; token body \texttt{\{\}}; username/email body \texttt{\{ identifier, password \}}.
\item \texttt{POST /auth/login} responses: \texttt{successCode: 200} for \texttt{token} with \texttt{access\_token}, \texttt{created}, \texttt{expired}; \texttt{successCode: 201} for \texttt{username} or \texttt{email} with \texttt{access\_token: \{ token, created, expired \}}; \texttt{BAD\_REQUEST} 400; \texttt{UNAUTHORIZED} 401; \texttt{CONFLICT} 409 (\texttt{DUPPLICATED\_USERNAME | DUPLICATED\_EMAIL | UNREGISTED\_EMAIL\_DOMAIN}); \texttt{UNPROCESSABLE} 422 (\texttt{INVALID\_USERNAME | INVALID\_EMAIL | INVALID\_PASSWORD}); \texttt{INTERNAL\_SERVER\_ERROR} 500 (\texttt{DATABASE\_ERROR | UNKNOWN\_ERROR}).
\item Error-code catalog: \texttt{UNKNOWN ERROR} 45000, \texttt{INVALID\_EMAIL} 45002, \texttt{INVALID\_USERNAME} 45003, \texttt{INVALID\_PASSWORD} 45004, \texttt{UNREGISTED\_EMAIL\_DOMAIN} 46001, \texttt{DUPLICATED\_EMAIL} 47001, \texttt{DUPLICATED\_USERNAME} 47002; backend \texttt{UNKNOWN ERROR} 2002 and \texttt{UNKNOWN DATABASE ERROR} 45000.
\item \texttt{GET /auth/validate} is listed as "Temporarily Disabled" with responses 204, 400, 409, 422, and 500; it is not the main login submission endpoint.
\item Session/Remember Me: \textbf{Remember Me} maintains the session across browser restarts (AC-02.3); expired sessions are denied on protected pages/APIs (FR-04.2), the inactivity timer resets on a qualifying request (FR-04.3), and expiration requires re-login (FR-04.4).
\item Data constraints: login accepts a registered username or university email address; CAPTCHA must be verified; backend checks activation status and restrictions; API account status values are inactive (0), preprofile (1), active (2), suspend (3).
\end{itemize}
\par\medskip

\hypertarget{sec:b11-source-alignment-notes}{}
\addcontentsline{toc}{subsubsection}{B11. Source Alignment Notes}
\subsubsection*{B11. Source Alignment Notes}
\par\medskip

\begin{itemize}
\item Inactivity timeout conflicts: FR-04.1 and US-04/AC-04.1 state 15 minutes, while NFR-SEC-7 states 30 minutes; \texttt{settings.md} recommends the stricter 15-minute value unless the product explicitly decides otherwise.
\item \texttt{API.md} defines \texttt{POST /auth/login} but no logout endpoint, although US-03 and FR-03.1 to FR-03.3 require logout, client cleanup, and backend session invalidation.
\item CAPTCHA is required by AC-02.4 and the use-case flow, but the documented \texttt{POST /auth/login} request body lists only \texttt{identifier} and \texttt{password} for username/email login.
\item Google login is required by FR-02.4 and AC-02.5 and modeled as E-02, but \texttt{API.md} does not define a dedicated Google OAuth login endpoint.
\end{itemize}
\par\medskip

\par\medskip
\hypertarget{sec:ui-prototype-specimens}{}
\addcontentsline{toc}{subsection}{UI Prototype Specimens}
\subsection*{UI Prototype Specimens}
\par\medskip
\begin{figure}[H]
\centering
\includegraphics[width=0.75\textwidth,height=0.42\textheight,keepaspectratio]{"../../prototype/login_desktop.png"}
\par\vspace{2pt}
{\footnotesize \textbf{UI Prototype:} Login - Desktop}
\end{figure}
\par\smallskip
\begin{figure}[H]
\centering
\includegraphics[width=0.75\textwidth,height=0.42\textheight,keepaspectratio]{"../../prototype/settings_login_desktop_destination.png"}
\par\vspace{2pt}
{\footnotesize \textbf{UI Prototype:} Settings: Login - Desktop destination}
\end{figure}
\par\smallskip

\end{document}
